% ==========================================================
% titles/title-11.tex -- Title XI: Amendment, Revision, and
% Constitutional Survival
% ==========================================================

\ConstTitle{Amendment, Revision, and Constitutional Survival}{title:amendment}

\Article{Ordinary Amendment}{art:ordinary-amend}

Provisions of this Constitution outside the foundational principles of Title I may be amended by a supermajority of four-fifths of the \gls{gls:sortition-legislature} (Article \ref{art:sortition-legislature}), followed by a national direct vote with sixty percent approval, followed by \gls{gls:constitutional-court} (Article \ref{art:const-court}) review confirming consistency with Title I principles.

\Article{The Failure Protocol}{art:failure}

\ArtSection{Acknowledgment}{art:failure:s-acknowledgment}
This Constitution is imperfect. Future generations will identify failures invisible to its authors. The measure of the designed system will not be the accuracy of its foresight but its ability to facilitate analysis of the material conditions and necessary self-correction.

\ArtSection{Extraordinary Constitutional Assembly}{art:failure:s-extraordinary-constitutional-assembly}
As established in Article \ref{art:const-assembly}, \ref{art:const-assembly:s-extraordinary-assembly}, if the \gls{gls:constitutional-court}, the \gls{gls:citizen-oversight-panels}, and the \gls{gls:transparency-engine} collectively and independently identify that the system is failing its foundational principles in ways beyond normal correction, they may jointly trigger an Extraordinary \gls{gls:constitutional-assembly} with full popular mandate to redesign the failing elements. See Article \ref{art:const-assembly}, \ref{art:const-assembly:s-extraordinary-assembly} for triggering conditions and procedures.

\ArtSection{The Anti-Sycophancy Clause}{art:failure:s-anti-sycophancy-clause}
No governance body and no \gls{gls:constitutional-assembly} may interpret this Constitution primarily to validate the decisions of the current governing generation. Interpretation is oriented toward the foundational principles and the interests of those who hold the least power, not toward institutional self-preservation.